\documentclass[11pt,a4paper]{article}
\usepackage{xeCJK}
\setCJKmainfont{Noto Serif CJK SC}
\setCJKsansfont{Noto Sans CJK SC}
\xeCJKsetup{CJKmath=true}
\usepackage{amsmath,amssymb}\usepackage{booktabs}\usepackage{array}
\usepackage[margin=2.3cm]{geometry}\usepackage{xcolor}\usepackage{titlesec}
\usepackage{enumitem}\usepackage{tikz}\usepackage{fancyhdr}
\definecolor{acc}{HTML}{25607F}\definecolor{pos}{HTML}{2E7D5B}\definecolor{neg}{HTML}{B0472E}
\definecolor{mut}{HTML}{69737D}\definecolor{rulec}{HTML}{C2C8CE}\definecolor{accs}{HTML}{E1EBF1}
\titleformat{\section}{\normalfont\Large\bfseries\color{acc}}{\thesection}{0.7em}{}
\titleformat{\subsection}{\normalfont\large\bfseries}{\thesubsection}{0.6em}{}
\setlist[itemize]{leftmargin=1.4em,itemsep=3pt,topsep=4pt}
\newcommand{\pp}[1]{\textcolor{pos}{#1}}\newcommand{\nq}[1]{\textcolor{neg}{#1}}
\pagestyle{fancy}\fancyhf{}\renewcommand{\headrulewidth}{0.4pt}
\lhead{\small\color{mut}半正定这道锁}\rhead{\small\color{mut}\thepage}
\setlength{\tabcolsep}{5pt}
\begin{document}

\begin{center}
{\Huge\bfseries 半正定这道锁}\\[7pt]
{\large\color{mut} Dota 2 Ability Draft · 胜负预测模型}\\[9pt]
{\small\color{mut} 训练数据 466{,}091 局(7.41d)\quad$\cdot$\quad 线上 \texttt{43.130.62.185/ad/} \quad$\cdot$\quad 2026-09-04}
\end{center}
\vspace{2pt}\noindent\rule{\textwidth}{1pt}\vspace{5pt}

\noindent 找了一年多的瓶颈,原来不在数据量、不在容量、也不在优化器 —— 是模型公式里一个没人注意的数学性质,它让模型在结构上说不出\textbf{“同类堆在一起反而更差”}。解开之后,一次拿到的提升超过此前所有改动的总和。

\section{结论}
\begin{center}\small
\begin{tabular}{@{}lrrrr@{}}
\toprule
模型 & 方向判对 & logloss & AUC & 配对 $t$\\
\midrule
改动前(线上原版) & 68.74\% & 0.58549 & 0.7546 & —\\
$+$ 队伍构成项 & 68.96\% & 0.58378 & 0.7565 & $+7.79$\\
$+$ 配合项切半 & 69.31\% & 0.57909 & 0.7609 & $+9.43$\\
\textbf{$+$ 跨队克制项} & \textbf{69.52\%} & \textbf{0.57733} & \textbf{0.7627} & \textbf{$+11.45$}\\
\bottomrule
\end{tabular}
\end{center}

三项各自独立显著(28 个实验下的 Bonferroni 阈值 $|t|>3.02$),可叠加。合计 logloss $-0.00816$,在 69{,}914 局留出测试集上多判对 \textbf{545 局}。

作为尺度参照:瞎猜的 logloss 是 0.693,改动前是 0.5855 —— 模型全部的本事是 0.1075。这一轮又挣了 0.0082,相当于\textbf{把已有功力再抬了 7.6\%}。

\section{那道锁}
原来的配合项是每件东西一个向量,两两点积:
\[ \text{座位配合} = \sum_{i<j}\langle e_i,e_j\rangle = \tfrac12\Bigl(\bigl\|\textstyle\sum e\bigr\|^2-\sum\|e\|^2\Bigr) \]

单个配对当然可以为负 —— 实测 204{,}480 个物品对里 \textbf{42.3\% 为负},最负 $-0.5077$(血魔本体 $\times$ 大圣归来)。所以“模型不能说两个东西互相拖累”是错的。

被锁死的是另一件事:$\|\sum e\|^2\ge 0$ 是硬约束,于是\textbf{一组彼此相似的东西,总配合必定为正}。

\vspace{4pt}
\noindent\fcolorbox{rulec}{white}{\parbox{0.97\textwidth}{\small
把同一件东西放 5 份(极端的“同类扎堆”):
\begin{center}
\begin{tabular}{@{}lll@{}}
对称 FM & $\mathrm{syn}=10\|e\|^2$ & 为负的 \nq{0 / 640} 件 —— 数学上不可能\\
允许负后 & & 为负的 \pp{165 / 640} 件,最负 $-2.452$\\
\end{tabular}
\end{center}
放 5 份是数学装置(AD 里每个技能只有一个),它测的是“这个方向上同类扎堆赚不赚”。对称写法下这个数是平方和,永远为正。}}

\vspace{8pt}
从 46.6 万局里挑出\textbf{同座位共现 $\ge$200 次、两种写法判断相反}的技能对,共 117 组:

\begin{center}\small
\begin{tabular}{@{}llrrr@{}}
\toprule
技能 A & 技能 B & 对称给 & 允许负后 & 共现\\
\midrule
死亡脉冲 & 自然之助 & \pp{$+0.314$} & \nq{$-0.028$} & 282\\
暗影波 & 自然之助 & \pp{$+0.252$} & \nq{$-0.063$} & 203\\
锚击 & 大圣归来 & \pp{$+0.385$} & \nq{$-0.020$} & 312\\
腐烂 & 竭心光环 & \pp{$+0.355$} & \nq{$-0.021$} & 942\\
神之谴戒 & 星破天惊 & \pp{$+0.263$} & \nq{$-0.026$} & 764\\
\bottomrule
\end{tabular}
\end{center}

前两组都是\textbf{两个范围治疗}放同一个座位 —— 血线该救得回来的一个就够。对称写法必须给正分,因为治疗类技能的向量彼此相近,而“相近”在点积下就是“高分”。

\subsection{一个三方对照}
月刃(弹射)、衰退光环(加攻击力)、灵能之刃(穿透)。两种写法的排序\textbf{完全相反},而实测站在后者一边:

\begin{center}\small
\begin{tabular}{@{}lrrrrr@{}}
\toprule
技能对 & 余弦 & 对称给 & 允许负后 & 共现 & 实测残差\\
\midrule
月刃 $\times$ 灵能之刃 & $+0.608$ & $+0.465$ & $+0.056$ & 162 & \nq{$-0.0471$}\\
月刃 $\times$ 衰退光环 & $+0.502$ & $+0.429$ & $+0.132$ & 231 & $-0.0051$\\
衰退光环 $\times$ 灵能之刃 & $+0.418$ & $+0.295$ & $+0.154$ & 297 & \pp{$+0.0394$}\\
\bottomrule
\end{tabular}
\end{center}

对称写法里“余弦”和“配合”是同一个数,两列必然同序 —— 它说月刃$+$灵能之刃最配。但这两个都是“让一次普攻打到更多目标”的机制,功能重叠;衰退光环是加攻击力,不在一个层面,能乘上去。\textbf{实测残差的方向和后者一致。}

\section{决定性对照}
把参数量和“是否允许负特征值”拆成两个维度,各跑一遍,验证集 logloss:

\begin{center}
\begin{tikzpicture}[font=\small]
\node[anchor=south] at (2.1,2.55) {\footnotesize\color{mut}只能加分(半正定)};
\node[anchor=south] at (7.1,2.55) {\footnotesize\color{mut}允许减分(不定)};
\node[rotate=90,anchor=south] at (-0.35,1.35) {\footnotesize\color{mut}61,440 参数};
\node[rotate=90,anchor=south] at (-0.35,-1.35) {\footnotesize\color{mut}122,880 参数};
\draw[rulec] (0.1,0.55) rectangle (4.1,2.35);
\node at (2.1,1.9) {\footnotesize\color{mut}对称 K=96(原版)};\node at (2.1,1.15) {\Large 0.5807};
\draw[acc,line width=1.2pt,fill=accs] (5.1,0.55) rectangle (9.1,2.35);
\node at (7.1,1.9) {\footnotesize\color{mut}切半 K=96};\node at (7.1,1.15) {\Large\color{acc} 0.5743};
\draw[rulec] (0.1,-2.35) rectangle (4.1,-0.55);
\node at (2.1,-1.0) {\footnotesize\color{mut}对称 K=192};\node at (2.1,-1.75) {\Large 0.5806};
\draw[rulec] (5.1,-2.35) rectangle (9.1,-0.55);
\node at (7.1,-1.0) {\footnotesize\color{mut}两张表 p/q d=96};\node at (7.1,-1.75) {\Large 0.5742};
\draw[acc,->,line width=1pt] (4.25,1.45) -- (4.95,1.45);
\node[above] at (4.6,1.5) {\scriptsize\color{acc}$-0.0064$};
\draw[rulec,->,line width=1pt] (2.1,0.4) -- (2.1,-0.4);
\node[right] at (2.2,0) {\scriptsize\color{mut}$-0.0001$};
\draw[rulec,->,line width=1pt] (7.1,0.4) -- (7.1,-0.4);
\node[right] at (7.2,0) {\scriptsize\color{mut}$-0.0001$};
\end{tikzpicture}
\end{center}

\noindent 参数翻倍在两种结构下都只值 0.0001;允许减分在两种参数量下都值 0.0064。\textbf{提升 100\% 来自结构。}

\vspace{6pt}
\noindent\fcolorbox{acc}{white}{\parbox{0.97\textwidth}{\small
\textbf{\color{acc}这条对照解释了之前所有的失败。} 自注意力、DeepSets、$K=64$、$K=192$、跨队自由参数 —— 之前一年多试过的东西,加的全是\textbf{容量},而缺的从来不是容量。每次都过拟合退回基线,现在知道为什么了。}}

\subsection{最终选的实现}
既然不是容量问题,就不用加参数。把 96 维切一半,后一半的点积取负:
\[ \text{配合}(i,j)=\sum_{k=1}^{48} e_{ik}e_{jk}\;\;\pp{-}\;\;\sum_{k=49}^{96} e_{ik}e_{jk} \]
参数 $640\times96=61{,}440$,一个不加;权重文件大小不变;前端改一行。符号是写死的不是学的,反传只多一个逐维乘法;全 $+1$ 时退化回原式,是同一份实现。

切半 $K{=}96$ 验证 0.5743,和参数翻倍的两张表版本 0.5742 打平,测试段还略好(0.5791 vs 0.5795)。

\section{另外两个模块}
“同类堆在一起会内耗”在三个层面各发生一次,而原模型三个都表达不了 —— 它是 10 个座位的线性相加,座位内又被半正定锁住。

\begin{center}\small
\begin{tabular}{@{}lllp{6.2cm}@{}}
\toprule
模块 & 层面 & 参数 & 实测缺口\\
\midrule
配合项切半 & 座位内 5 件 & 0(不加) & 见上\\
队伍构成项 & 队伍内 5 座位 & 9 & 4 个打钱位比 1 个低 \nq{6.7pp},$t=-13.9$\\
跨队克制项 & 我方 vs 敌方 & 9,600 & 整体超零模型 $+10.5\%$,$z=+14.4$\\
\bottomrule
\end{tabular}
\end{center}

\textbf{队伍构成项}是一条 8 段分段线性,按队伍打钱分总和查表。实测形状是倒 U —— 一个打钱位都没有也吃亏($-0.0188$),1$\sim$2 个最好,4 个 $-0.0650$。上线后这条残差从 $t=-5.36$ 压到 $t=-1.47$。

\textbf{跨队克制项}用 15 条冻结轴描述敌方阵容(力/敏/智/近战本体数 $+$ 打钱/治疗/爆发浓度 $+$ 8 条队伍向量主成分),每件东西学 15 个标量表示“对这类阵容有多在意”。写成差的形式,反对称由结构保证。

\section{克制:证据扎实,吃到的有限}
\begin{center}\small
\begin{tabular}{@{}lrrl@{}}
\toprule
案例 & 低档 & 高档 & 检验\\
\midrule
死亡旋风 vs 敌方力量英雄数 & \nq{$-0.0101$} & \pp{$+0.0260$} & 300 个同频段技能安慰剂,$z=+3.70$\\
全能大 vs 敌方打钱位数 & \nq{$-0.0153$} & \pp{$+0.0530$} & $t=+5.77$\\
冰魂本体 vs 敌方治疗浓度 & \nq{$-0.0105$} & \pp{$+0.0115$} & $t=+4.20$\\
\bottomrule
\end{tabular}
\end{center}

三个都过了\textbf{池内混杂控制}:固定“这一局一共有几个力量英雄”,只看“其中几个在对面”。三个总数档全部单调、方向一致,合并后 $t$ 分别是 $+3.76$ / $+6.51$ / $+4.32$。这排除了“这种局本身就算不准”的解释。

\vspace{4pt}
\noindent\fcolorbox{acc}{white}{\parbox{0.97\textwidth}{\small
\textbf{\color{acc}AD 特有的一个陷阱。} 技能池由这局的 10 个英雄决定,所以\textbf{同英雄的技能必然同池} —— 跨队共现中位数 15{,}665 局,是不同英雄(1{,}450)的 \textbf{10.8 倍}。这会让任何系统性偏差因为功效过高而“显著”。做了两个阴性对照:冰魂本体的效应在按大招位置分层后仍是 $+0.0216$(不分层 $+0.0220$);割裂 vs 焦渴(同为血魔技能,共现 21{,}102 局)残差 $+0.0019$、$t=0.64$,没有伪信号。}}

\subsection{没解决的:旋风只吃到 1/12}
模型\textbf{方向学对了} —— 力量轴是死亡旋风 15 条轴里最大的一条($+0.0079$,是全体标准差的 2.82 倍),智力轴是负的($-0.0052$),和实测的对偶关系吻合。但幅度:
\begin{center}\small
\begin{tabular}{@{}ll@{}}
模型给的 & 敌方每多一个力量英雄 $\to$ 胜率 $+0.20$ 个百分点\\
实测需要 & 胜率 $+2.5$ 个百分点\qquad\textbf{差 12 倍}\\
\end{tabular}
\end{center}

病根是\textbf{统一收缩}。7{,}680 个系数里绝大多数应该是零,训练用统一学习率 $+$ 统一早停,停点由“大多数是噪声”决定,真信号被和噪声一起叫停。试了三条路都没解决:
\begin{itemize}
\item \textbf{换参数化}(技能对 $\to$ 聚合轴,样本量 $1{,}450\to43{,}547$):系数 0.0077,没动
\item \textbf{换轴集合}(12 条 $\to$ 15 条):$0.0077\to0.0079$,没动
\item \textbf{放松收缩}:学习率 $\times4$ 崩掉;显式 L2 让系数翻 2.1 倍,但聚合 logloss 变差 0.0007,且旋风\textbf{在全体里的相对位置没变}($2.82\to2.70$)
\end{itemize}

最后一条最说明问题:解耦 weight decay 对所有参数施加同一个 $\lambda$,和早停一样是统一的;而 Adam 的自适应缩放恰好把“样本量差异”抹平了。真正的“按证据分配”需要绕开 Adam 重写求解器。

没继续做,因为预期收益已经很小:那条 3$\sim$5 档只有 $n=1001$、标准误 0.0144,实测的 $+0.0503$ 本身 95\% 区间是 $[+0.022,+0.079]$。\textbf{模型给的保守值未必是错的。}

\section{过程里的教训}
\begin{itemize}
\item \textbf{随机初始化的组件不能和已收敛的组件联合训练。} 犯了两遍:克制项用残差矩阵初始化而主模型从零,把模型卡在 0.587;换配合项时辅助头形状变了只能随机初始化,把热启动的 0.5806 拽到 0.6070。都得靠“先冻结、分段训”解决。
\item \textbf{从收敛解微调必须大幅降学习率。} $8\times10^{-4}$ 是从零训练的值,起点在最优解上时一轮就把参数踢飞。降到 $10^{-4}$ 仍然不如从零重训。
\item \textbf{float32 下的梯度自检会误报。} 梯度被 batch 归一化压到 $10^{-8}$ 量级,float32 的损失只有 $\sim10^{-7}$ 精度,差分测到的全是浮点噪声。必须隔离出那一项、在 float64 下查 —— 有一次因此误判“梯度全错”,实际公式没问题。
\item \textbf{早停要真的加上。} 训练脚本一直没有耐心早停,$K=192$ 那次最优在第 11 轮却跑到第 19 轮,堵了后面四个实验一个多小时。
\item \textbf{前端展示要跟着模型口径改。} 换成切半后拆解面板还在用旧公式,同一座位卡片显示 0.22、展开显示 0.890;总分行只算座位分不含新加的两项,导致 $\mathrm{sigmoid}(\text{左}-\text{右})$ 和显示的胜率对不上(44.3\% vs 53.2\%)。都是玩家一眼能看出的自相矛盾。
\item \textbf{浏览器缓存要根治。} 静态文件加 mtime 版本号,否则各人加载的 JS 版本不一致 —— 这也是之前“房主看不到、别人能看到”的根因。
\end{itemize}

\vfill
\noindent\rule{\textwidth}{0.4pt}\\
{\footnotesize\color{mut}
线上模型 \texttt{model\_plan15\_clr0.25}:配合项切半 $K{=}96$(前 48 维加分 / 后 48 维减分)$+$ 队伍构成 8 段 $+$ 跨队克制 15 轴。
权重 443 KB 随页面下发,浏览器本地计算,服务端不参与推理。
上线前 JS 与 Python 逐样例对拍 logit 最大差 $6.85\times10^{-8}$,反对称自检 $8.88\times10^{-16}$。
回滚包 \texttt{/root/ad/rollback/before\_3mod\_20260903\_200805.tgz}。\\
网页版:\texttt{https://claude.ai/code/artifact/2b59df25-ba79-4e81-a860-69eb8c48daa0}}
\end{document}
